\chapter{PART VII · 面向未来}


\newpage

未来不是未来才来的。

它已经在敲门了。AI的能力每六个月就在某些维度上翻一番。AI独角兽企业的估值在以无法解释的速度攀升。AI原生组织正在以我们从未见过的商业模式挑战传统巨头。

而大多数组织，还在用工业时代的思维运营知识经济时代的员工，还在用20世纪的管理逻辑应对21世纪的竞争现实。

我们讨论了AI时代的思维重构、流程重构、组织重构、技能重构，以及变革的阻力维度。但所有这些讨论，最终要回到一个问题：

接下来怎么办？

三个章节，对应三个面向未来的视角：
\begin{itemize}
\item 德鲁克会如何写《动荡时代的管理》——他在1980年写了这本书，如果他在2026年重写，他会改变什么核心观点
\item 管理者的下一步：具体行动框架——把本书讨论的框架转化为可操作的行动步骤
\item 结语：AI只是工具，人才是目的——回到最根本的问题

\end{itemize}
这不是一本关于AI的书。这是一本关于AI时代如何做组织的书。技术是背景，人是主体，逻辑是主线。
